\homeworktitle{Домашняя работа}{Рациональные уравнения с модулем и параметром}
\studentline

\begin{routebox}
  \textbf{Ключевой факт:}
  \[
    \frac{A(x)}{B(x)}=0
    \quad\Longleftrightarrow\quad
    \begin{cases}
      A(x)=0,\\
      B(x)\ne 0.
    \end{cases}
  \]
  \textbf{Порядок работы:}
  \begin{enumerate}[label=\arabic*.,leftmargin=1.8em,itemsep=0.2em,topsep=0.25em]
    \item исследуйте нули числителя;
    \item найдите условие нужного количества нулей;
    \item проверьте каждый найденный ноль в знаменателе;
    \item объедините условия и аккуратно запишите ответ.
  \end{enumerate}
\end{routebox}

\section*{Часть I. Разминка}

\begin{homeworktask}[Ноль дроби]{1}
  Решите уравнение
  \[
    \frac{(x-2)(x+3)}{x-2}=0.
  \]
  Коротко объясните, почему один из нулей числителя не является корнем
  исходного уравнения.
\end{homeworktask}
\answerspace{14mm}

\Needspace{18\baselineskip}
\begin{homeworktask}[Количество корней числителя]{2}
  Определите количество корней уравнения
  \[
    2|x|-x=b
  \]
  в зависимости от значения параметра \(b\).

  Сначала представьте левую часть по промежуткам, затем заполните таблицу.
  \[
    2|x|-x=
    \begin{cases}
      \rule{27mm}{0.4pt}, & x\ge 0,\\[2mm]
      \rule{27mm}{0.4pt}, & x<0.
    \end{cases}
  \]

  \begin{center}
    \renewcommand{\arraystretch}{1.4}
    \begin{tabularx}{0.78\textwidth}{>{\centering\arraybackslash}p{0.3\textwidth}>{\centering\arraybackslash}X}
      \toprule
      \rowcolor{TableHeader}
      Значение \(b\) & Количество корней \\
      \midrule
      \(b<0\) & \rule{30mm}{0.4pt} \\
      \(b=0\) & \rule{30mm}{0.4pt} \\
      \(b>0\) & \rule{30mm}{0.4pt} \\
      \bottomrule
    \end{tabularx}
  \end{center}
\end{homeworktask}

\Needspace{18\baselineskip}
\section*{Часть II. Обязательные задачи}

\begin{homeworktask}[Близкий аналог разобранной задачи]{3}
  Найдите все значения параметра \(a\), при каждом из которых уравнение
  \[
    \frac{|2x|-x-1-a}{x^2-x-a}=0
  \]
  имеет ровно два различных корня.

  \textbf{В решении должны быть явно указаны:}
  \begin{itemize}[itemsep=0.15em,topsep=0.25em]
    \item условия, при которых числитель имеет два различных нуля;
    \item оба нуля числителя;
    \item результаты их подстановки в знаменатель;
    \item окончательный ответ по параметру.
  \end{itemize}
\end{homeworktask}
\shortanswerline

\newpage
\begin{homeworktask}[Самостоятельное применение алгоритма]{4}
  Найдите все значения параметра \(a\), при каждом из которых уравнение
  \[
    \frac{3|x|-x-2-a}{x^2-x-a}=0
  \]
  имеет ровно два различных корня.
\end{homeworktask}
\solutionmap

\Needspace{16\baselineskip}
\section*{Часть III. Осмысление и перенос метода}

\begin{homeworktask}[Найдите ошибку]{5}
  Ученик исследовал исходное уравнение
  \[
    \frac{|3x|-2x-2-a}{x^2-2x-a}=0
  \]
  и написал:

  \begin{quote}
    «Числитель имеет два различных нуля при \(a>-2\). Следовательно,
    ответ: \(a>-2\)».
  \end{quote}

  Объясните ошибку ученика и запишите правильный ответ.
\end{homeworktask}
\answerspace{18mm}

\newpage
\begin{homeworktask}[Задание со звёздочкой: параметр внутри модуля]{6*}
  Найдите все значения параметра \(a\), при каждом из которых уравнение
  \[
    \frac{2|x-a|-x-2}{x^2-1}=0
  \]
  имеет ровно два различных корня.

  \begin{hintbox}
    Рассмотрите случаи \(x\ge a\) и \(x<a\). Здесь параметр перемещает
    точку излома графика \(y=2|x-a|\) вдоль оси \(Ox\), поэтому готовый
    шаблон основной задачи нельзя применять механически.
  \end{hintbox}
\end{homeworktask}
\solutionmap

\begin{selfcheck}
  \begin{itemize}[label=$\square$,leftmargin=1.7em,itemsep=0.25em,topsep=0.2em]
    \item Я сохранил условия ветвей при раскрытии модуля.
    \item Я проверил каждый ноль числителя в знаменателе.
    \item Я различаю значения переменной \(x\) и значения параметра \(a\).
    \item В ответе указаны промежуток и все исключаемые точки.
  \end{itemize}
\end{selfcheck}
